% Migrated from proposal/proposal.tex (5 October 2026) on 14 October 2026.
\chapter{Methodology}
\label{ch:methodology}

\section{Draft Status and Study Design}
The study plan below is transferred from the proposal dated 5 October 2026.
It describes intended work; it does not report a completed study.
\todo{Revise this plan against the supervisor meeting of 13 October 2026:
remove category quotas; use English primarily with German as fallback;
include apps without a privacy policy; extend the keyword space to mental
health trackers; exclude connected-device analysis; specify two-coder
procedures and justify the sampling and access rules. The ordering and
remaining thresholds require a documented decision.}

The study will use structured content analysis~\parencite{krippendorff2019content}
with an adapted cognitive walkthrough. Figure~\ref{fig:pipeline} summarises
the workflow. Each positive code will link to a verbatim quotation, source
location and capture date. Sampling rules, reference trees, the codebook and
prompts will be versioned and frozen as a dated release before the main study.
Model identifiers, settings and raw outputs will be retained so that reported
counts can be traced to evidence.

\section{Literature Review} Searches will cover privacy harms
\parencite{solove2006taxonomy,calo2011boundaries}, tracker privacy, and
methods for policy analysis~\parencite{wilson2016creation,zimmeck2019maps}
and LLM-assisted annotation. The log will record databases, queries, dates and
inclusion decisions. A comparison matrix will distinguish each study's domain,
sources and treatment of harms.

\section{Pilot} A five-app pilot (three fitness, two fertility) will test
collection and coding boundaries and estimate capture time, coding units,
manual effort and model cost. Pilot apps will be recoded with the final
instrument and remain in the sample, but will not enter the validation set.

\begin{figure}[t]
\centering
\resizebox{\textwidth}{!}{%
\begin{tikzpicture}[
  font=\scriptsize,
  bx/.style={draw,rounded corners=2pt,align=center,text width=29mm,
             minimum height=13mm,inner sep=2.5pt},
  arr/.style={-{Stealth[length=1.6mm]},semithick}]
  \node[bx] (a) at (-7.2,0) {\textbf{Pilot and sample}\\47 ranked store apps\\3 preinstalled apps};
  \node[bx] (b) at (-3.6,0) {\textbf{Capture}\\documents for 50 apps\\walkthroughs for 12};
  \node[bx] (c) at (0,0) {\textbf{Code}\\coverage, stance\\protection claims};
  \node[bx] (d) at (3.6,0) {\textbf{Check}\\blind manual coding\\audit missed mentions};
  \node[bx] (e) at (7.2,0) {\textbf{Compare}\\RQ1--RQ4\\with source evidence};
  \draw[arr] (a)--(b); \draw[arr] (b)--(c);
  \draw[arr] (c)--(d); \draw[arr] (d)--(e);
\end{tikzpicture}}
\caption{Planned workflow. The pilot will refine the instruments before they
are frozen. Later changes will be logged and affected material recoded.}
\Description{Five boxes connected by arrows: pilot and sample, capture,
code, check, compare. Documents cover 50 apps and walkthroughs cover 12.}
\label{fig:pipeline}
\end{figure}

\section{App Sampling}
\paragraph{Candidates.} A fixed list of 19 keywords will cover both domains in
German and common English listing terms, such as \emph{Schrittzähler},
\emph{Workout}, \emph{Periode}, \emph{Eisprung} and \emph{period tracker}. Each
keyword will be searched in the German App Store and the German Google Play
store with the collectors in Appendix~\ref{app:scripts}, and the first 20
results per keyword and store will be recorded with the date of collection.
Listings found by several keywords will be merged.

\paragraph{Ranking by exposure.} The study concerns the communication that
users meet, so candidates will be ranked by how widely an app is used, not by
how well it is rated. Neither store publishes download counts
(Table~\ref{tab:storefront-fields}), so the number of ratings will stand in for
use. Sunyaev et al.\ selected the 300 most frequently rated apps per store on
the same grounds and reported a Spearman correlation of 0.89 between rating
count and download count on Android~\parencite{sunyaev2015availability}. The
relationship will be checked again on the candidate pool: Google Play shows a
worldwide rating count and a worldwide install figure for each app, and their
Spearman rank correlation will be reported. A pilot collection on 22 October
2026 with all 19 keywords found 159 Google Play apps with both figures; their
correlation was 0.96, and 0.94 for the 135 apps with at least 1{,}000 ratings.
The App Store publishes no downloads, so the check rests on Google Play. Eligible listings need at least
1{,}000 ratings. Each store will be ranked separately, because App Store counts
describe the German storefront and Google Play counts the worldwide store, so
the two counts cannot be compared directly.

\paragraph{Store balance.} Both stores will have equal weight, as in earlier
reviews that took the same number of apps from each
store~\parencite{sunyaev2015availability,oloughlin2019reviewing}. Apps will be
taken alternately from the two ranked lists, starting with Google Play, whose
operating system runs on most German smartphones~\parencite{statcounter2026vendor}.
An app already selected through the other store will be skipped, so each app
counts once and both of its listings are captured. Selection stops at 47 apps,
about half from each store.

\paragraph{Preinstalled apps.} Three places are reserved for health apps that
ship with the phone. They reach users without a download and so cannot be
ranked by the same measure. The rule covers phone vendors with at least 10\,\%
of German mobile use; in September 2026 these were Samsung (39\,\%), Apple
(28\,\%) and Xiaomi (21\,\%)~\parencite{statcounter2026vendor}. Their
preinstalled health or fitness trackers are Apple Health and Apple Fitness on
the iPhone and Samsung Health on Samsung phones. No preinstalled tracker on
Xiaomi's international software could be identified, so Xiaomi contributes
none, and Huawei and Google fall below the threshold with about 1.5\,\% each. A
preinstalled app that also appears in the ranked lists moves to this group, and
the ranked selection continues to 47. Apple Health has no store listing, so it
will be analysed through its documentation and the device.

Neither store publishes German figures for every listing field, and the two
stores split their fields differently. Apple shows one summary rating per
country or region, tied to the storefront where the customer bought the
app~\parencite{apple2026ratingsoverview}. Google Play has shown phone users
ratings from their registered country since November
2021~\parencite{grinsted2021making}. Download counts are reported only to an
app's developer, in App Store Connect~\parencite{apple2026acquisition} and in the
Play Console~\parencite{google2026playstatistics}. These sources do not say
which public field follows which rule, so Table~\ref{tab:storefront-fields}
rests on a field check. On 22 October 2026 it requested the same listings from
the German, US and British storefronts and compared the values from Apple's
search and lookup endpoint~\parencite{apple2022itunessearch}, the App Store
product pages, and the Google Play product pages read with an open-source
parser~\parencite{jomingyu2024googleplayscraper}. On the App Store, every rating
figure describes the requested storefront and no worldwide total is published.
On Google Play, the rating count and the install figures were the same in every
storefront, while the average rating and the one-to-five-star breakdown
changed: the breakdown is the country's distribution scaled to the worldwide
count. Both stores show an app's position in a category chart of the
storefront. The check covered few apps in one session, so it shows how the
fields behave and does not enter the sample. The two stores therefore count
ratings on different bases.

\begin{table}[t]
\centering
\begin{singlespace}
\caption{Public listing data by geographic basis. Observed on 22 October 2026
for up to six App Store and five Google Play listings in the German (DE), US
and British (GB) storefronts; fields shown only on product pages were checked
for one app per store in DE and US. \emph{Market-specific}: the value changed
with the storefront. \emph{Worldwide}: it was identical across storefronts.
A field check, not part of the sample.}
\label{tab:storefront-fields}
\small
\setlength{\aboverulesep}{0.25ex}\setlength{\belowrulesep}{0.35ex}
\begin{tabularx}{\textwidth}{@{}>{\raggedright\arraybackslash}p{2.3cm}>{\raggedright\arraybackslash}X>{\raggedright\arraybackslash}X@{}}
\toprule
\textbf{Field} & \textbf{App Store} & \textbf{Google Play} \\
\midrule
Rating count & \textbf{Market-specific}, also for the current version; no
  worldwide total. Yazio: 446{,}768 (DE), 51{,}379 (US), 12{,}351 (GB)
  & \textbf{Worldwide}, labelled \enquote{reviews} on the page; no country
  count. Home Workout: 4{,}045{,}258 (DE, US), 4{,}045{,}260 (GB) \\
\midrule
Average rating & \textbf{Market-specific}, also for the current version.
  Yazio: 4.62 (DE), 4.70 (US)
  & \textbf{Market-specific}. MyFitnessPal: 3.94 (DE), 4.41 (US) \\
\midrule
One-to-five-star breakdown & \textbf{Market-specific}, on the product page
  only. Yazio five-star: 349{,}509 (DE), 42{,}759 (US)
  & \textbf{Mixed}: the country's distribution scaled to the worldwide count.
  MyFitnessPal one-star: 428{,}497 (DE), 205{,}443 (US) \\
\midrule
Written reviews & A few review texts on the product page; \textbf{no public
  count}
  & Undocumented embedded count, not displayed, \textbf{market-specific}.
  MyFitnessPal: 32{,}560 (DE), 454{,}136 (US) \\
\midrule
Downloads or installs & \textbf{Not public}; reported to the developer
  only~\parencite{apple2026acquisition}
  & \textbf{Worldwide} band on the page (\enquote{100M+}) and an
  undocumented exact count (Home Workout: 303{,}142{,}208 in all three);
  country figures to the developer only~\parencite{google2026playstatistics} \\
\midrule
Chart position & \textbf{Market-specific} badge on the product page (Yazio:
  no.~3 in Health \& Fitness, DE) and public category chart pages
  & \textbf{Market-specific} badge when ranked (MyFitnessPal: no.~2 top
  grossing in Health \& Fitness, US; none in DE) \\
\midrule
Search result order & Endpoint order per storefront; not verified as the
  native store rank
  & Page order per storefront; not verified as the native store rank \\
\bottomrule
\end{tabularx}
\end{singlespace}
\end{table}

\paragraph{Sample size and eligibility.} The target is \textbf{50 apps}: 47
ranked store apps and 3 preinstalled apps. The number of preinstalled apps
follows from the vendor rule, and the ranking fills the rest of the target.
There are no quotas per domain; fitness and fertility apps compete in one
ranking, and the resulting mix will be reported. Apps must collect health or
body-related tracking data, be free to install, offer German or English text
and have a reachable policy. Device companions remain eligible if they record
or display tracking data. Duplicates and DiGA (Germany's reimbursable digital
health applications) will be excluded. Apps serving both domains will be
flagged. Unreachable policies will be logged in the screening results before
exclusion from text analysis. If fewer than 47 eligible store apps are found,
all will be retained and the shortfall reported. Thresholds will not be lowered
after inspecting results.

\section{Material and Language} For each app, the core corpus will include
the policy, store descriptions, structured store disclosures, the main product
page, and privacy/security pages linked directly from it or the policy. A
fixed supplementary search will capture provider-hosted privacy blog posts
from the preceding 12 months. Ads encountered during collection will form an
opportunistic supplement. The absence of ads or blog posts will not count as
absence of a harm in the core corpus.

Language is recorded per source. German text will be primary, with English
used when German is unavailable; results describe this German-first corpus.
If both exist, differences will be logged. Original wording will be coded.
Each source will be archived with URL, date, language, file hash and app
version where applicable. Shared policies will be coded
once and linked to all relevant apps.

\section{App Walkthroughs} Twelve distinct apps will be inspected, six per domain
with three Android and three iOS apps in each. Selection will follow the
ranking, alternating Android and iOS slots without reusing an app; the balance
serves coverage, not platform comparison. Apps requiring an unavailable paired
device, payment or identity verification will be skipped with a recorded
reason, and unfilled slots reported. Each session will use a fresh install,
synthetic data and a study account where required.

Tasks will cover finding collected data, inferred data, breach consequences
and incident responses, reaching the policy, stopping optional third-party
sharing, exporting data and deleting the account. At each step the researcher
will ask Wharton et al.'s four questions: would the user pursue the intended
goal, notice the action, connect it to that goal and recognise progress
\parencite{wharton1994cognitive}? The pilot will fix starting points and
time limits. Outcomes will distinguish completion, information not found,
unsupported tasks and access barriers. Screen recordings will preserve the
text encountered. This is an expert inspection, not a test of real users.

\section{Reference Harm Trees} Before main coding, a tree will be built for
each domain from prior tracker research. For fitness, sources are a survey of
tracker inferences, threats and vulnerabilities~\parencite{niksirat2024wearable}
and users' harm scenarios~\parencite{gabriele2020understanding}. For fertility,
sources are analyses of policies, data flows and interfaces
\parencite{shipp2020period,malki2024exploring,hudig2025intimate} and post-Roe
user concerns~\parencite{cao2024deleted}. Sources found in the logged search
will be added with a recorded reason. Each pathway will name the data, potential
adversary, access conditions, capability and possible consequence. Required
conditions will be connected explicitly, following Knabenhans et al.'s
AND/OR approach. A prediction made for the user will not itself establish
third-party access (Appendix~\ref{app:example}). As in Knabenhans et al., capabilities include incapabilities, such as a user
being unable to use a feature without agreeing to data sharing.

The component origins will distinguish core functionality, digital context
and implementation choices within the stated model. For every app and
pathway, an applicability record will note the required feature, its
evidence (listing, policy or walkthrough), the access and context
assumptions, the decision (applicable, not applicable or uncertain) and the
reason. A harm category applies to an app when at least one of its pathways
does. The trees are reference models of possible harms, not complete
descriptions of the apps' actual systems. Uncertain cases will be counted and
reported, not silently dropped, so that sparsely documented providers do not
vanish from the gap analysis.

\section{Harm Categories} The starting categories are Knabenhans et al.'s
physical, economic, psychological, autonomy/deprivation, discrimination,
reputational and societal harms~\parencite[App.~A]{knabenhans2026harms},
drawn from Citron and Solove and the CyberPeace Institute. Following them, a
harm is an impairment of the ability to function as one otherwise would in
one's usual context~\parencite[\S4.1]{knabenhans2026harms}; here, the baseline
is the user's situation without the modelled pathway. A relationship tag will
retain Citron and Solove's relationship harms~\parencite{citron2021harms},
which the list omits. New harms found in the pilot will be documented rather
than forced into a category.

\section{Coding Units and Axes} To code a text is to attach labels that
record what it says. A unit will be one paragraph or one coherent block of
interface text. Adjacent text will be available for context, but supporting quotations must
be identified precisely. The codebook will define each label, a fitting
example and a counter-example. Table~\ref{tab:axes} summarises the three axes.
Appendix~\ref{app:example} shows the procedure on a constructed example.

\begin{table}[t]
\centering
\begin{singlespace}
\caption{Coding axes and boundaries. All example wording is constructed.
Several labels may apply to one unit.}
\label{tab:axes}
\small
\setlength{\aboverulesep}{0.25ex}\setlength{\belowrulesep}{0.35ex}
\begin{tabularx}{\textwidth}{@{}L{1.6cm}L{2.6cm}>{\raggedright\arraybackslash}X>{\raggedright\arraybackslash}X@{}}
\toprule
\textbf{Axis} & \textbf{Label} & \textbf{What it records} & \textbf{Example or boundary} \\
\midrule
A: Coverage & Component & Data collected, stored or disclosed, or a system requirement
  & \enquote{We store your cycle dates.} \\
\rowrule
 & Capability & Distinguishing users, tracking/profiling, inferring an attribute, or an incapability
  & \enquote{We estimate your fertile window.} \\
\rowrule
 & Concrete harm & A named negative consequence and its category
  & \enquote{Disclosure could cause distress.} \\
\rowrule
 & Generic risk & Possible misuse without a concrete consequence
  & \enquote{Your data could be used against you.} \\
\rowrule
 & No relevant content & None of the coverage labels applies
  & A contact address alone \\
\rowrule
 & Additional details & Breach consequences, incident response and relationship impacts
  & Code concrete response steps separately from a general legal-notice promise \\
\midrule
B: Stance & Acknowledged & A named harm presented as possible
  & \enquote{A disclosure could cause distress.} \\
\rowrule
 & Qualified & Uncertainty or likelihood language, with its target recorded
  & \enquote{A leak is unlikely, but could cause distress.} \\
\rowrule
 & Minimised & A named harm dismissed without supporting reasons in the unit or adjacent text
  & \enquote{Such distress is no cause for concern.} \\
\rowrule
 & Responsibility shifted & Responsibility for a named harm placed on the user
  & \enquote{You alone must prevent harassment caused by sharing your location.} \\
\rowrule
 & Denied/prevented & A named harm presented as impossible or prevented
  & \enquote{This makes harassment impossible.} \\
\midrule
C: Protection & Measure and scope & Claimed measure, actor, data, outcome and conditions
  & On-device storage is a measure. A no-sale pledge is a practice commitment. \\
\rowrule
 & Framing and strength & Data-flow or harm framing, with partial, scoped or absolute protection
  & \enquote{This reduces exposure} differs from \enquote{No one can harm you.} \\
\rowrule
 & Reassurance & Broad safety language without a specific measure
  & \enquote{Your data is safe with us.} \\
\rowrule
 & Potential mismatch & Human-reviewed tension between the claim and a pathway within its scope
  & A residual risk outside the claim's scope is not sufficient \\
\bottomrule
\end{tabularx}
\end{singlespace}
\end{table}

\subsection{Axis A: Coverage} All explicitly mentioned nodes will be
recorded, including whether each is asserted, presented as possible or denied.
The highest level reached will be a separate summary, so a harm
mention does not erase accompanying data-flow information. Feature descriptions can mention inference without presenting it as a risk.

\subsection{Axis B: Stance} Stance labels will be linked to the
particular harm and supporting span, rather than assigned to the paragraph
as a whole. A statement may acknowledge a harm and qualify its likelihood,
and the coder will distinguish uncertainty about a breach from uncertainty
about its consequences. Minimisation will be judged from wording in the unit
or adjacent text, never from inferred intent; qualification alone does not
establish it. RQ2 thus covers stance toward named harms; dismissive language
without one is coded as reassurance. These distinctions adapt Zou et
al.~\parencite{zou2019might} and will be refined in the pilot.

\subsection{Axis C: Protection} The coder will record the measure
and claimed benefit, with actor, data and conditions where specified. Any
technical measure claimed to limit exposure or misuse counts, PET-labelled or
not, including ordinary security controls; non-technical commitments are kept
separate. A \emph{potential mismatch} requires a pathway within the claim's
scope and an explanation of why the measure misses its necessary conditions;
partial protection is not rejected merely because other risks remain. Vague
scope will be recorded as unclear, and no deployed system judged safe or
unsafe.

\section{LLM-Assisted Coding} Polisis and recent LLM approaches are precedents for
data-practice analysis~\parencite{harkous2018polisis,rodriguez2024llm}. The
harm-focused scheme will use prompts and a model selected on the pilot, then
frozen~\parencite{tornberg2024best}; an open-weight model will be preferred if
it performs close to a hosted one, keeping the run repeatable. Prompts will
work in stages: relevant spans, then stance for detected harms, then
protection details. All retained harm mentions and potential mismatches will
be reviewed manually, so the model acts as a screener whose prompts favour
recall. The author will decide tree applicability and mismatches, logging
evidence and reasoning; uncertain outputs or missing spans enter manual
review rather than becoming negative codes. Each unit will run three times;
agreement is a stability check, not validation, and labels returned once are
reviewed rather than discarded. Raw output is kept apart from adjudicated
labels, so performance is not measured after correction.

\section{Validation and Missed Mentions} Ten apps, five per domain, will be
drawn at random with a recorded seed from providers and documents not used in
the pilot, stratified by the language of the coded policy where both occur;
the numbers of independent providers and unique documents will be reported.
The author will code them fully before seeing model output, and they will not
be used to tune prompts. Because prompt wording alone can change
conclusions~\parencite{baumann2025llm}, two paraphrased prompts will also be
run and headline counts compared. Disagreements will be adjudicated with each
label's source hidden; two apps will be recoded after at least two weeks,
before drafting, to estimate intra-coder consistency. A second coder will be
sought for applicability, minimisation and mismatch decisions; failing that,
the supervisor will review a sample, including judged non-mismatches, with
disagreements recorded before discussion. A detection is correct only with
the right label, unit and evidence span, and the staged pipeline is scored end
to end, so first-stage misses count. Reported will be per-code precision
(share of model labels supported), recall (share of manual mentions found)
and model--human Cohen's $\kappa$~\parencite{cohen1960coefficient}, with
counts, prevalence and Gwet's AC1, since rare codes can depress $\kappa$
despite high agreement~\parencite{feinstein1990high,gwet2008computing}.
App-level detection is also reported, by domain and, where counts allow, by
source and language. Exploratory 95\,\% intervals will resample providers to
respect shared sources; with ten apps they will be wide.

A separate missed-mention audit, reported apart from validation, will search
the remaining corpus for harm-related vocabulary in both languages and check
every unlabelled hit manually. To expose the vocabulary's blind spots, 100
units outside the hits (or all, if fewer), allocated across domains, languages
and sources in advance, will also be read. Recall will be reported as
unestimable where manual positives are too few. For any code with fewer than
20 validation positives or recall below 0.90, all relevant material will be
checked manually or corpus-wide absence claims withheld; these are review
triggers, not proof of accuracy.

\section{Analysis and Reporting} Results will be descriptive and reported
separately for fitness and fertility apps. For RQ1, each harm's denominator
will include only applicable apps with the required core sources captured. A
source the provider does not offer is recorded as such; only sources that
exist but cannot be retrieved reduce the denominator. Inapplicable pathways
and text with no identified mention will have different labels. Counts will be per app and surface, with
shared-provider results also shown after grouping repeated providers.
RQ2 will report harm--stance combinations and their denominators. RQ3 will
summarise measures, scope, claim strength and manually reviewed mismatches.
Examples that support and challenge the expected gap will both be retained.

For RQ4, the principal policy--interface comparison will use the same 12
inspected apps; the 50-app document analysis gives the broader overview, and
other surface pairs report how many apps have both surfaces. At the fixed
decision points, onboarding before consent and the privacy-settings screen,
text shown directly, behind a visible link or not found within the task limit
will be distinguished. Inconsistent claims will be reported apart from mere
differences in detail. Absence will mean \enquote{no mention was identified
in the captured material}, with the instrument's measured limitations stated.

\section{Ethics and Feasibility} No human participants are planned.
Walkthroughs use synthetic data, store terms will be checked before automated
retrieval, and any review requirements confirmed before collection. The
pilot's estimates will fix a minimum feasible study (core documents, the blind
validation set, the walkthroughs) before collection; if time runs short,
supplementary marketing collection is reduced first and other changes are
agreed with the supervisor.

\section{Limitations} The sample favours free, frequently rated apps with
reachable policies and does not represent all German users' exposure. Rating
counts are a proxy for use: Google Play's count is worldwide rather than
German, and apps that prompt users to rate collect more ratings. The corpus is a
dated snapshot, walkthroughs cover selected paths, and reference trees depend
on the chosen literature. Unless a second coder joins, one person provides the
reference labels, so model agreement is not independent human-coder
reliability~\parencite{mcdonald2019reliability}. Missed mentions may remain.
The study cannot establish provider intent, actual harms or effects on users.
